\section*{Chapitre \ref{ch:g11:mutations} --- Mutations et variation génétique}

\begin{solution}{exo:g11:mutations:1}
Un changement de la séquence de \omterm{def:g10:universal-dna:nucleotide}{nucléotides} de l'\omterm{def:g10:universal-dna:information}{ADN}, transmis aux
descendants de la \omterm{def:g10:cells-common-unit:cell}{cellule}. Substitution, insertion, délétion.
\end{solution}

\begin{solution}{exo:g11:mutations:2}
\texttt{ATGCCAGAAT}~: substitution (T en A à la position 6).
\texttt{ATGCCTGAT}~: délétion d'un A. \texttt{ATGCCTTGAAT}~: insertion
d'un T.
\end{solution}

\begin{solution}{exo:g11:mutations:3}
Somatique~: dans une \omterm{def:g10:cells-common-unit:cell}{cellule} du corps, transmise aux seuls descendants
de cette \omterm{def:g10:cells-common-unit:cell}{cellule} chez l'individu. Germinale~: dans un gamète ou son
précurseur, transmise à la descendance et à toutes ses \omterm{def:g10:cells-common-unit:cell}{cellules}. Seules
les \omterm{def:g11:mutations:mutation}{mutations} germinales entrent dans l'hérédité.
\end{solution}

\begin{solution}{exo:g11:mutations:4}
Le rayonnement ultraviolet, les rayons X, le benzopyrène de la fumée de
tabac (mais aussi la radioactivité, l'aflatoxine). L'ultraviolet soude
ensemble deux thymines voisines d'un brin, ce qui déforme l'hélice.
\end{solution}

\begin{solution}{exo:g11:mutations:5}
Le brin complémentaire intact~: comme chaque brin détermine l'autre, la
séquence correcte du segment abîmé peut être reconstruite à partir de
son partenaire.
\end{solution}

\begin{solution}{exo:g11:mutations:6}
Normales~: environ 6~\%~; déficientes en réparation~: 0.01~\%. La
réparation améliore la survie d'un facteur d'environ 600 à cette dose.
\end{solution}

\begin{solution}{exo:g11:mutations:7}
$20/10^9 = \num{2e-8}$. Non~: la boîte ne révèle que les \omterm{def:g10:cells-common-unit:cell}{cellules} déjà
résistantes~; savoir si l'antibiotique a causé la \omterm{def:g11:mutations:mutation}{mutation} demande le
test des répliques.
\end{solution}

\begin{solution}{exo:g11:mutations:8}
Chaque réplique reçoit des \omterm{def:g10:cells-common-unit:cell}{cellules} des mêmes colonies d'origine. Si
l'antibiotique provoquait la résistance au contact, les \omterm{def:g10:cells-common-unit:cell}{cellules}
touchées seraient quelques-unes au hasard sur chaque réplique, à des
positions sans rapport. Trouver la résistance toujours aux mêmes
positions signifie que ces colonies d'origine, poussées sans le
médicament, contenaient déjà des \omterm{def:g10:cells-common-unit:cell}{cellules} résistantes.
\end{solution}

\begin{solution}{exo:g11:mutations:9}
Erreurs par copie du \omterm{def:g10:universal-dna:gene}{gène}~: $2 \times 3000 \times 10^{-9} =
\num{6e-6}$ (les deux brins). Parmi les $10^{12}$ \omterm{def:g10:cells-common-unit:cell}{cellules} produites à
la dernière division, environ $\num{6e6}$ portent une \omterm{def:g11:mutations:mutation}{mutation} nouvelle
dans ce \omterm{def:g10:universal-dna:gene}{gène}.
\end{solution}

\begin{solution}{exo:g11:mutations:10}
Une \omterm{def:g11:mutations:mutation}{mutation} somatique dans une \omterm{def:g10:cells-common-unit:cell}{cellule} de bourgeon a donné sa couleur
à la branche~; les boutures sont ce tissu et la gardent donc. Les
\omterm{def:g10:cells-common-unit:cell}{cellules} germinales des fleurs dérivent de la même branche, mais un
changement dans un \omterm{def:g10:universal-dna:gene}{allèle} n'est transmis qu'à la moitié des gamètes, et
les ovules ont été fécondés par un pollen porteur de l'\omterm{def:g10:universal-dna:gene}{allèle}
ordinaire~; le nouvel \omterm{def:g10:universal-dna:gene}{allèle}, s'il est récessif, reste caché chez les
plantules.
\end{solution}

\begin{solution}{exo:g11:mutations:11}
L'\omterm{def:g10:cell-metabolism:metabolism}{enzyme} manquante répare les dommages dus aux ultraviolets~; les
ultraviolets ne pénètrent que la peau. L'intestin ne reçoit aucun
ultraviolet, ne subit aucune lésion de ce genre, et n'a pas besoin de
cette réparation.
\end{solution}

\begin{solution}{exo:g11:mutations:12}
Les lésions sont proportionnelles à la dose~: un dixième des
ultraviolets fait un dixième des lésions par \omterm{def:g10:cells-common-unit:cell}{cellule} et par heure ---
soit encore des milliers. Le coup de soleil est la mort de nombreuses
\omterm{def:g10:cells-common-unit:cell}{cellules}~; en dessous de cette dose, les \omterm{def:g10:cells-common-unit:cell}{cellules} survivent, mais les
lésions qui échappent à la réparation deviennent des \omterm{def:g11:mutations:mutation}{mutations}. Pas de
brûlure signifie que les dommages étaient survivables, non qu'ils
étaient absents.
\end{solution}

\begin{solution}{exo:g11:mutations:13}
Ni l'un ni l'autre en soi. Dans une région où sévit le paludisme,
l'exemplaire unique élève la survie, et l'\omterm{def:g10:universal-dna:gene}{allèle} reste fréquent malgré
la maladie des porteurs doubles~; là où le paludisme est absent, il ne
fait que provoquer la maladie et il est rare. C'est l'environnement, et
le nombre d'exemplaires, qui décident de la valeur d'une \omterm{def:g11:mutations:mutation}{mutation}.
\end{solution}

\begin{solution}{exo:g11:mutations:14}
Les \omterm{def:g10:cells-common-unit:cell}{cellules} sensibles meurent~; les résistantes, déjà présentes, se
divisent librement et en une journée la culture est entièrement
résistante. Une cure interrompue trop tôt laisse en vie les derniers
survivants --- les plus résistants --- pour repeupler un corps
désormais enrichi en résistance~; mener la cure à son terme les tue
avant qu'ils ne se multiplient.
\end{solution}

\begin{solution}{exo:g11:mutations:15}
La plupart des \omterm{def:g11:mutations:mutation}{mutations} sont neutres (hors des \omterm{def:g10:universal-dna:gene}{gènes} ou sans effet sur
la \omterm{def:g10:chemistry-of-life:families}{protéine}), une minorité est nuisible, un très petit nombre est
utile~; étant aléatoires, elles surviennent sans égard à leur utilité,
mais la sélection garde les utiles et élimine les nuisibles. L'utilité
dépend de l'environnement, si bien qu'une \omterm{def:g11:mutations:mutation}{mutation} nuisible aujourd'hui
peut être celle qui compte quand les conditions changent. Sur des
milliers de générations, avec des milliards d'individus, même des
\omterm{def:g11:mutations:mutation}{mutations} utiles rares apparaissent à répétition~: l'évolution est
bâtie sur le reliquat trié, non sur la \omterm{def:g11:mutations:mutation}{mutation} moyenne.
\end{solution}

\begin{solution}{pb:g11:mutations:1}
\textbf{1.} Pour que chaque colonie soit un clone d'origine connue, et
pour que les colonies soient assez séparées pour être distinguées sur
les répliques.

\textbf{2.} Quelques \omterm{def:g10:cells-common-unit:cell}{cellules} du sommet de chaque colonie, dans la même
disposition géométrique que sur la boîte, parce que le tampon est
pressé à plat sans glisser.

\textbf{3.} Que les \omterm{def:g10:cells-common-unit:cell}{cellules} résistantes étaient présentes dans ces
deux colonies d'origine~: trois répliques indépendantes ne s'accordent
pas par hasard.

\textbf{4.} Des colonies résistantes à des positions quelconques et
sans rapport sur les trois boîtes, puisque les \omterm{def:g10:cells-common-unit:cell}{cellules} induites ne
seraient pas les mêmes d'une réplique à l'autre.

\textbf{5.} La \omterm{def:g11:mutations:mutation}{mutation} avant contact~: la résistance se rapporte aux
colonies d'origine, qui n'ont jamais rencontré le médicament.

\textbf{6.} Les \omterm{def:g10:cells-common-unit:cell}{cellules} des deux positions résistantes poussent sur
streptomycine~; celles des autres positions, non.

\textbf{7.} Une partie seulement, probablement~: la \omterm{def:g11:mutations:mutation}{mutation} est
survenue à une division quelconque pendant la croissance de la colonie,
et seuls les descendants de cette \omterm{def:g10:cells-common-unit:cell}{cellule} sont résistants. Pour
trancher, il faudrait étaler un nombre connu de \omterm{def:g10:cells-common-unit:cell}{cellules} de la colonie
sur streptomycine et compter les survivantes.

\textbf{8.} À la 20\textsuperscript{e} des 23 divisions~:
$1/2^{20-1}$\dots plus simplement, la \omterm{def:g10:cells-common-unit:cell}{cellule} mutante se divise ensuite
3 fois de plus~: $2^3 = 8$ \omterm{def:g10:cells-common-unit:cell}{cellules} résistantes sur $2^{23}$, soit
environ une sur un million. À la 3\textsuperscript{e} division~: l'une
des 8 \omterm{def:g10:cells-common-unit:cell}{cellules} alors présentes, soit un huitième de la colonie.

\textbf{9.} Même quelques \omterm{def:g10:cells-common-unit:cell}{cellules} résistantes parmi les $10^4$ environ
que le velours prélève suffisent à fonder une colonie sur la boîte à
antibiotique~; le test détecte la présence de \omterm{def:g10:cells-common-unit:cell}{cellules} résistantes, non
leur proportion.

\textbf{10.} L'argument repose sur la position~: ce n'est que si la
géométrie est conservée qu'une colonie d'une réplique peut être
identifiée à une colonie de la boîte d'origine.

\textbf{11.} $30/\num{10000} = \num{3e-3}$.

\textbf{12.} Environ $\num{3e-3}$ \omterm{def:g11:mutations:mutation}{mutation} pour $10^7$ divisions~:
$\num{3e-10}$ par division.

\textbf{13.} Du même ordre que le taux d'erreur sur un \omterm{def:g10:universal-dna:nucleotide}{nucléotide}~: la
\omterm{def:g11:mutations:mutation}{mutation} de résistance est essentiellement une substitution précise qui
échappe à la réparation.

\textbf{14.} $10^9$ divisions à $\num{3e-10}$~: en moyenne une fraction
de \omterm{def:g11:mutations:mutation}{mutation}, si bien que la plupart des cultures n'en contiennent
aucune ou quelques \omterm{def:g10:cells-common-unit:cell}{cellules} résistantes seulement (chaque \omterm{def:g11:mutations:mutation}{mutation}
étant ensuite multipliée par les divisions qui la suivent).

\textbf{15.} Une \omterm{def:g11:mutations:mutation}{mutation} survenue à une division précoce est héritée
par une grande fraction de la culture, et le comptage est donc élevé~;
une \omterm{def:g11:mutations:mutation}{mutation} tardive ne donne que quelques \omterm{def:g10:cells-common-unit:cell}{cellules} résistantes. Comme
le moment est aléatoire, des cultures répétées donnent des comptages
très dispersés --- ce qui est en soi un signe que les \omterm{def:g11:mutations:mutation}{mutations}
apparaissent avant l'antibiotique, et non en réponse à lui.

\textbf{16.} Rares~: les \omterm{def:g11:mutations:mutation}{mutations} surviennent peu souvent par
division. Aléatoires~: elles surviennent indépendamment de leur
utilité. L'expérience établit le caractère aléatoire (l'argument des
positions)~; les comptages de la partie III mesurent la rareté.

\textbf{17.} Il a trié la population~: il a tué les \omterm{def:g10:cells-common-unit:cell}{cellules} sensibles
et laissé les résistantes se multiplier --- c'est la sélection.

\textbf{18.} La \omterm{def:g11:mutations:mutation}{mutation} a fourni l'\omterm{def:g10:universal-dna:gene}{allèle} de résistance, au hasard,
avant tout besoin~; l'antibiotique a trié la population, en n'en
gardant que les porteurs.

\textbf{19.} La \omterm{def:g11:mutations:mutation}{mutation} change une molécule impliquée dans l'action du
médicament, non la croissance ni l'aspect de la colonie en l'absence de
médicament~; une différence génétique n'a d'effet visible que dans un
environnement où elle compte.

\textbf{20.} Les positions identiques sur les trois répliques ont
prouvé que les \omterm{def:g11:mutations:mutation}{mutations} de résistance apparaissent avant l'exposition
à l'antibiotique et indépendamment d'elle~; le taux mesuré vaut environ
$\num{3e-10}$ par division cellulaire.
\end{solution}
